\documentclass[10pt,a4paper]{amsart}
\usepackage[margin=19mm]{geometry}
\usepackage{amsmath,amssymb,mathtools}
\usepackage[scr=rsfs,cal=euler]{mathalfa}
\usepackage{xfrac}
\usepackage[unicode,hidelinks,bookmarks=false]{hyperref}
\newcommand{\ie}{i.e.\ }
\newcommand{\cf}{cf.\ }
\newcommand{\resp}{respectively\ }
\newcommand{\Subsection}{\subsection}
\providecommand{\nobreakdash}{\nobreak\hbox{-}}
\setlength{\emergencystretch}{2em}
\multlinegap0pt
\usepackage{amssymb}
\newtheorem{lemma}{Lemma}
\newcommand{\R}{\mathbb R}
\newcommand{\consubsp}{\mathcal{I}} % Notation for space spanned by a column vector of ones
\newcommand{\eqrowsum}{\mathcal{M}} % Notation for space containing all matrices of the form 1c'
\newcommand{\gensn}[1]{|#1|} % Notation for general consensus seminorm (on matrices)
\newcommand{\restnorm}[1]{\|#1\|_R} % Notation for matrix norm on restricted (quotient) space
\newcommand{\vecsn}[1]{|#1|} % Notation for general consensus norm (on column vectors)
\title{\LARGE \bf On Determining the Convergence Rate of an Infinite Product of Stochastic Matrices*}
\author{Ron Ofir, A. Stephen Morse}
\date{Source: arXiv:2607.03623v1 (CC BY 4.0)}
\begin{document}
\maketitle

\noindent\emph{Source and licence.} \LARGE \bf On Determining the Convergence Rate of an Infinite Product of Stochastic Matrices*. Version arXiv:2607.03623v1 (CC BY 4.0), distributed by arXiv under the Creative Commons Attribution 4.0 International licence, http://creativecommons.org/licenses/by/4.0/, as recorded in the arXiv licence metadata for this exact version and shown on the abstract page https://arxiv.org/abs/2607.03623v1. Full licence notice: \texttt{licenses/arxiv-2607.03623-CC-BY-4.0.txt}.\par\smallskip

\noindent\textit{Editorial scope.} Counted unit: the article's lemma lem:restnorm (source0.tex lines 179--187). The article's complete original proof of it is delivered with it (source0.tex lines 189--216, one \texttt{proof} environment of 28 lines). No mathematics is written, repaired or re-derived: the statement and the proof are the article's own lines, in the article's own notation, and no retained line is substituted or re-flowed.  No further source-local context is needed: every cross-reference of every retained line resolves inside the two delivered spans.  Nothing was omitted from any retained span, so the package carries no omission record.  No retained line contains a citation, so the wrapper carries no bibliography and no bibliography entry.  No retained line contains a figure environment, an image-inclusion command or drawing code, so no image is delivered and the package declares no asset dependency.  Editorial formatting only, in addition: an AMS \texttt{amsart} preamble that restates the article's own declarations and macros used by the retained lines, namely the environments \texttt{lemma} on separate counters, together with the standard packages those lines need.  The retained environments print in this document as Lemma 1 in order of appearance, whereas the article prints the same items with the numbering of its own full document.  The complete native source of the article as distributed in the arXiv e-print for this exact version is bundled unchanged as \texttt{sources/204404/source0.tex}.
\begin{lemma}\label{lem:restnorm}
    Let~$P$ be a full-rank matrix with~$\ker Q = \consubsp$, and let~$\gensn{\cdot} : \eqrowsum^{n \times n} \to \R$ be a consensus seminorm. There exists a norm~$\restnorm{\cdot} : \R^{(n-1) \times (n-1)} \to \R$ such that
    \begin{equation}
        \gensn{M} = \restnorm{\bar M}
    \end{equation}
    for every~$M \in \eqrowsum^{n \times n}$, where~$\bar M \in \R^{(n-1) \times (n-1)}$ is the unique matrix for which~$QM = \bar M Q$.

    Furthermore, if~$\gensn{\cdot}$ is submultiplicative, then so is~$\restnorm{\cdot}$.
\end{lemma}

\begin{proof}
    We will show that the function~$\restnorm{\cdot} : \R^{n-1 \times n -1} \to \R$ given by
    \begin{equation}\label{eq:restnorm}
        \restnorm{A} = \vecsn{M},
    \end{equation}
    where~$M \in \eqrowsum^{n \times n}$ is such that~$PM = AP$, is well-defined. In particular, it will be shown that for every~$A \in \R^{(n-1) \times (n-1)}$ there exist infinitely many~$M \in \eqrowsum^{n \times n}$ such that~$PM = AP$, but that the right hand side of~\eqref{eq:restnorm} is the same for every such~$M$.
    
    Since~$P$ has full row rank, it has a right inverse~$P^{-1} \in \R^{n \times n-1}$. In addition, since~$\ker P = \consubsp$, for every~$A \in \R^{n-1 \times n-1}$,~$P^{-1}AP\mathbf{1} = 0$ so~$P^{-1}AP \in \eqrowsum^{n \times n}$, and~$P(P^{-1}AP) = AP$. Fix~$A \in \R^{n-1 \times n-1}$, and suppose~$M_1,M_2 \in \eqrowsum^{n\times n}$ are such that
    \begin{equation}
        AP = PM_1 = PM_2.
    \end{equation}
    Then, since~$\ker P = \consubsp$, we have that~$M_1-M_2 = \mathbf{1} c$ for some~$c \in \R^{1 \times n}$. It follows from the shift-invariance of consensus seminorms that~$\vecsn{M_1} = \vecsn{M_2}$, so the right hand side of~\eqref{eq:restnorm} is the same for any choice of~$M$.

    It remains to show that~$\restnorm{\cdot}$ is indeed a norm. Let~$A_1,A_2 \in \R^{n-1 \times n-1}$ and~$M_1,M_2 \in \eqrowsum^{n \times n}$ such that~$PM_i=A_iP, i=1,2$. Let~$\alpha \in \R$. Then,~$\alpha(A_1 + A_2)P = \alpha P(M_1 + M_2)$, so
    \begin{equation}
        \restnorm{\alpha(A_1 + A_2)} = \vecsn{\alpha(M_1 + M_2)} \le |\alpha|(\vecsn{M_1} + \vecsn{M_2}),
    \end{equation}
    and this proves homogeneity and the triangle inequality. Nonnegativity of~$\restnorm{\cdot}$ is immediate from the nonnegativity of seminorms. It remains only to show that~$\restnorm{A} = 0$ if and only if~$A = 0$. Since~$P \cdot 0 = 0 \cdot P$,~$\restnorm{0} = 0$. Suppose now that~$A \in \R^{n-1 \times n-1}$ is such that~$\restnorm{A} = 0$. This implies that~$AP = P\mathbf{1}c$ for some~$c \in \R^{1 \times n}$. But, since~$P\mathbf{1}c = 0$, we have that~$APP^{-1} = P\mathbf{1}cP^{-1} = 0$, so~$A = 0$.
    
    Finally, suppose that~$\gensn{\cdot}$ is submultiplicative. Let~$A_1,A_2 \in \R^{n-1 \times n-1}$ and~$M_1,M_2 \in \eqrowsum^{n \times n}$ such that~$PM_i=A_iP, i=1,2$. Then,
    \begin{equation}
        (A_1 A_2)P = A_1 PM_2 = PM_1 M_2,
    \end{equation}
    so
    \begin{equation}
        \restnorm{A_1 A_2} = \vecsn{M_1 M_2} \le \vecsn{M_1}\vecsn{M_2} = \restnorm{A_1}\restnorm{A_2}.
    \end{equation}
\end{proof}

\end{document}
